\documentclass[11pt]{article}
\usepackage[margin=0.75in]{geometry}
\usepackage{ctex}
\usepackage{longtable}
\usepackage{array}
\usepackage{booktabs}
\renewcommand{\arraystretch}{1.25}

\begin{document}

\begin{center}
{\Large AP宏观经济学互动Lab简表}
\end{center}

\begin{longtable}{p{0.13\textwidth} p{0.25\textwidth} p{0.56\textwidth}}
\toprule
\textbf{Unit} & \textbf{Lab} & \textbf{Lab描述} \\
\midrule
\endfirsthead
\toprule
\textbf{Unit} & \textbf{Lab} & \textbf{Lab描述} \\
\midrule
\endhead

Unit 1 & Scarcity and Choice & 通过资源分配观察稀缺性、选择和机会成本。 \\
Unit 1 & PPC and Opportunity Cost & 用PPC图判断效率、不可达点、经济增长和机会成本。 \\
Unit 1 & Comparative Advantage & 比较两个经济体的机会成本，判断比较优势和贸易收益。 \\
Unit 1 & Economic Systems & 对比市场、命令和混合经济中资源配置方式的差异。 \\

Unit 2 & Measuring Output & 用GDP组成、名义GDP和实际GDP理解总产出衡量。 \\
Unit 2 & Measuring Prices & 用CPI、GDP Deflator和通胀率衡量价格水平变化。 \\
Unit 2 & Labor Market Indicators & 区分就业、失业和非劳动力，并计算失业率与劳动参与率。 \\
Unit 2 & Business Cycle & 识别扩张、峰值、衰退和谷底，并连接产出、失业和通胀变化。 \\

Unit 3 & AD Components & 调节消费、投资、政府支出和净出口，观察总需求变化。 \\
Unit 3 & AD-AS Equilibrium & 用AD、SRAS和LRAS分析宏观均衡、价格水平和实际产出。 \\
Unit 3 & Output Gaps & 判断衰退缺口和通胀缺口，并选择合适的政策方向。 \\
Unit 3 & Fiscal Policy Simulator & 观察政府支出、税收和转移支付如何通过乘数影响AD。 \\
Unit 3 & Multiplier Lab & 用MPC计算支出乘数和税收乘数，理解政策冲击大小。 \\

Unit 4 & Bank Balance Sheet & 用T-account理解存款、准备金、超额准备金和贷款。 \\
Unit 4 & Money Multiplier & 观察准备金率和初始存款如何影响最大货币创造。 \\
Unit 4 & Reserve Market & 用准备金供给和准备金需求分析准备金利率变化。 \\
Unit 4 & Money Market & 用货币供给和货币需求解释名义利率变化。 \\
Unit 4 & Loanable Funds & 用可贷资金市场分析储蓄、投资和实际利率。 \\
Unit 4 & Monetary Policy & 观察货币政策如何影响利率、投资、AD和短期产出。 \\

Unit 5 & Phillips Curve & 用短期和长期Phillips Curve理解通胀与失业的关系。 \\
Unit 5 & Expectations Adjustment & 观察通胀预期变化如何移动短期Phillips Curve。 \\
Unit 5 & Economic Growth & 用PPC和LRAS外移表示长期经济增长。 \\
Unit 5 & Deficits and Crowding Out & 用可贷资金市场分析赤字、债务和挤出效应。 \\
Unit 5 & Money Growth and Inflation & 观察货币增长在长期中如何影响价格水平和通胀。 \\

Unit 6 & Balance of Payments & 区分经常账户和金融账户，理解国际收支关系。 \\
Unit 6 & Foreign Exchange Market & 用外汇市场分析汇率、升值、贬值和净出口。 \\
Unit 6 & Capital Flows & 观察利率差异如何影响资本流动和汇率。 \\
Unit 6 & Policy in an Open Economy & 分析财政或货币政策在开放经济中对汇率和净出口的影响。 \\

\bottomrule
\end{longtable}

\end{document}
